\section*{Problemas}

\subsection*{Enunciados de los problemas}

% Recordar
\begin{problema}
	\label{prob:97c4388}
	Enuncie, sin realizar ningún cálculo, la función de densidad $f(x)$ y el soporte de una variable $X \sim U(a,b)$, y escriba de memoria las fórmulas de $\mathbb{E}[X]$ y $\text{Var}(X)$.
\end{problema}

% Comprender
\begin{problema}
	\label{prob:9f6e77c}
	El tiempo de espera de un autobús en una parada se modela como $X \sim U(0, 15)$ minutos. Un pasajero afirma: ``como la exponencial no tiene memoria, la uniforme tampoco, así que si ya esperé 10 minutos, la probabilidad de esperar 5 minutos más es la misma que al principio''. Explique por qué esta afirmación es falsa para la distribución uniforme, calculando explícitamente $P(X \ge 15 \mid X \ge 10)$ y comparándola con $P(X \ge 5)$.
\end{problema}

% Aplicar
\begin{problema}
	\label{prob:8fcc221}
	Con el mismo tiempo de espera $X \sim U(0, 15)$ minutos:
	\begin{enumerate}
		\item Calcule la probabilidad de que el tiempo de espera sea menor a 5 minutos.
		\item Calcule el cuantil $q_{0.75}$ (el tiempo bajo el cual espera el 75\% de los pasajeros).
	\end{enumerate}
\end{problema}

% Analizar
\begin{problema}
	\label{prob:5794e09}
	Sea $X$ una variable continua con soporte acotado en $[a,b]$. Demuestre, usando multiplicadores de Lagrange sobre la restricción $\int_a^b f(x)\,dx=1$, que la distribución $U(a,b)$ maximiza la entropía diferencial $h(X) = -\int f(x)\log f(x)\,dx$ entre todas las distribuciones con ese soporte.
\end{problema}

% Evaluar
\begin{problema}
	\label{prob:126be41}
	Un ingeniero de simulación necesita generar muestras de $Y \sim \text{Exp}(\lambda=2)$ y afirma: ``el método de la transformada inversa ($Y = F^{-1}(U)$ con $U \sim U(0,1)$) solo funciona si $F$ tiene una forma cerrada explícita, así que no puedo usarlo para distribuciones sin CDF invertible en forma cerrada''. Evalúe esta afirmación: ¿es correcta la condición que exige el método? Justifique con la demostración de por qué $Y=F^{-1}(U)$ tiene CDF $F$, y aclare qué es lo que realmente se necesita para aplicar el método.
\end{problema}

% Crear
\begin{problema}
	\label{prob:c1c6ece}
	Diseñe un ejemplo original (contexto y parámetros $a,b$) que se modele naturalmente con una distribución Uniforme continua, distinto de los ejemplos de tiempos de espera ya vistos en esta sección. Plantee la pregunta de probabilidad de interés y calcule $\mathbb{E}[X]$ y $\text{Var}(X)$ para los parámetros que eligió.
\end{problema}

\subsection*{Soluciones detalladas}

% Recordar
\begin{solproblema}[prob:97c4388]
	La función de densidad es $f(x) = \dfrac{1}{b-a}$ para $x \in [a,b]$ y $0$ en otro caso. Los momentos son $\mathbb{E}[X] = \dfrac{a+b}{2}$ y $\text{Var}(X) = \dfrac{(b-a)^2}{12}$.
\end{solproblema}

% Comprender
\begin{solproblema}[prob:9f6e77c]
	La afirmación es \textbf{falsa}: la falta de memoria es una propiedad exclusiva de la distribución exponencial (y geométrica en el caso discreto), no se transfiere a otras distribuciones. Con $X \sim U(0,15)$, la CDF es $F(x)=x/15$, así que:
	\begin{align*}
		P(X \ge 15 \mid X \ge 10) = \dfrac{P(X\ge15)}{P(X\ge10)} = \dfrac{0/15}{5/15} = 0,
	\end{align*}
	mientras que $P(X\ge5) = 10/15 = 2/3 \approx 0.667$. Como $0 \neq 2/3$, la propiedad de falta de memoria claramente no se cumple: al haber esperado ya 10 minutos, el pasajero sabe con certeza que el autobús llegará en los siguientes 5 minutos (el soporte acotado en 15 se lo garantiza), justo lo opuesto a lo que predeciría la falsa analogía con la exponencial.
\end{solproblema}

% Aplicar
\begin{solproblema}[prob:8fcc221]
	Con $F(x) = x/15$:
	\begin{enumerate}
		\item $P(X<5) = F(5) = 5/15 = 1/3 \approx 0.333$.
		\item $q_{0.75}$ satisface $q_{0.75}/15 = 0.75$, es decir $q_{0.75} = 11.25$ minutos.
	\end{enumerate}
\end{solproblema}

% Analizar
\begin{solproblema}[prob:5794e09]
	Maximizamos $h[f] = -\int_a^b f(x)\log f(x)\,dx$ sujeto a $\int_a^b f(x)\,dx=1$. El Lagrangiano es $\mathcal{L}[f] = -\int_a^b f(x)\log f(x)\,dx - \lambda\left(\int_a^b f(x)\,dx - 1\right)$. La derivada funcional respecto a $f$ da $-\log f(x) - 1 - \lambda = 0$, de donde $f(x) = e^{-(1+\lambda)}$ es constante en $[a,b]$. La normalización exige $f(x) = 1/(b-a)$, que es exactamente la densidad de $U(a,b)$. Por lo tanto, entre todas las distribuciones con soporte en $[a,b]$, la uniforme es la que maximiza la entropía diferencial — la ``menos informativa'', consistente con no imponer ninguna preferencia dentro del intervalo.
\end{solproblema}

% Evaluar
\begin{solproblema}[prob:126be41]
	La afirmación es \textbf{incorrecta}. La demostración de que $Y=F^{-1}(U)$ tiene CDF $F$ es general:
	\begin{align*}
		P(Y\le y) = P(F^{-1}(U)\le y) = P(U\le F(y)) = F(y),
	\end{align*}
	usando solo la monotonía de $F$ (y por tanto de $F^{-1}$) y que $U\sim U(0,1)$. Lo único que realmente se necesita es que $F$ sea invertible (estrictamente creciente en su soporte) — no que $F^{-1}$ tenga una \emph{forma cerrada explícita}. Cuando $F^{-1}$ no tiene forma cerrada, el método sigue siendo válido en principio; simplemente se evalúa $F^{-1}$ numéricamente (por ejemplo, por bisección o algún solver). Para $\text{Exp}(\lambda=2)$, de hecho, sí existe forma cerrada: $F(y)=1-e^{-2y} \implies Y = -\tfrac{1}{2}\ln(1-U)$, así que ni siquiera aplica la supuesta limitación en este caso concreto.
\end{solproblema}

% Crear
\begin{solproblema}[prob:c1c6ece]
	\textbf{Ejemplo:} un generador de números aleatorios de hardware produce voltajes de ruido térmico distribuidos uniformemente en $X \sim U(-5, 5)$ milivoltios. Pregunta de interés: ¿cuál es la probabilidad de que el voltaje medido exceda $3$ mV en valor absoluto, $P(|X|>3)$? Los momentos son $\mathbb{E}[X] = \dfrac{-5+5}{2} = 0$ y $\text{Var}(X) = \dfrac{(5-(-5))^2}{12} = \dfrac{100}{12} \approx 8.33$.
\end{solproblema}
